\section*{Assembly cut after the complete singleton KL payment}

All logarithms are base two.  For a genuine Boolean function put
\(p(S)=\widehat f(S)^2\), \(\pi(A)=\sum_{S\supseteq A}p(S)\), and retain
the exact signed restriction identity
\[
 \Delta_iU=B_i-\Gamma_i+\sum_{A\subseteq[n]\setminus\{i\}}K_{i,A}.
 \tag{1}
\]
The KL row \(K_{i,A}\) uses the signed upper-set correlation; no absolute
value or arbitrary-law relaxation is made.

\subsection*{Frozen candidate and exact accepted inputs}

Freeze
\[
 R_{4,3/4}(f)=G(f)-4I(f)-\frac34U(f).
 \tag{2}
\]
The sharp root theorem and the complete singleton-row theorem state
\[
 K^{(0)}\le {2\over\ln2}M_1,
 \qquad
 K^{(1)}\le {1\over2}V_1+{2\over\ln2}W_2,
 \tag{3}
\]
where
\[
 M_1=\sum_i\min\{p(\varnothing),p(\{i\})\},\quad
 W_2=\sum_{|S|=2}p(S),
 \tag{4}
\]
\[
 V_1=\sum_{i\ne j}
 \bigl(p(\{j\})+p(\{i,j\})\bigr)
 \log_2{\pi(\{j\})\over p(\{j\})+p(\{i,j\})}.
 \tag{5}
\]
The source proof of the second inequality is valid: in each physical child
\(g=f_{i,\pm}\), Booleanity gives
\(|\widehat g(\{j\})|\le\operatorname{Inf}_j(g)\).  This is exactly the
line that controls the signed posterior denominator.  It is not available
for arbitrary spectral laws.

Let
\[
 D_{\ge2}=\sum_i(B_i-\Gamma_i)
       +\sum_i\sum_{|A|\ge2}K_{i,A}.
 \tag{6}
\]
Since \(M_1+W_2\le I\), substitution of (3) into (1), summed in \(i\),
gives the all-dimensional theorem
\[
 \boxed{G\le D_{\ge2}+{1\over2}V_1+{2\over\ln2}I.}
 \tag{7}
\]
Consequently it is sufficient, only on tensor-prime Boolean functions with
\(R_{4,3/4}>0\), to prove
\[
 \boxed{
 D_{\ge2}+{1\over2}V_1
 \le\left(4-{2\over\ln2}\right)I+{3\over4}U.}
 \tag{GU-cut}
\]
This implies \(G\le4I+3U/4\), then \(U\le16I\), and the audited source
reduction gives \(H\le34I\).  Constants and zero-influence functions are
handled separately.

\subsection*{Minimal-counterexample structure}

If (2) is positive somewhere, choose a counterexample of minimum dimension
and pass to a positive-residual tensor-prime factor.  Every proper physical
child has nonpositive residual.  Therefore every relevant coordinate obeys
\[
 \Delta_iG-\frac34\Delta_iU
 \ge4\operatorname{Inf}_i(f)+R_{4,3/4}(f).
 \tag{8}
\]
No extremizer in unbounded dimension is assumed.  Equation (8) is a valid
structural hypothesis for attacking (GU-cut), but it does not by itself pay
either term on its left.

\subsection*{Two transfers that are not available}

First, \(V_1/2\le U\) is false for genuine Boolean spectra.  On four bits,
take the truth table (lexicographic binary input order)
\[
 (1,-1,1,1,1,1,-1,1,-1,-1,-1,1,1,-1,-1,-1).
 \tag{9}
\]
Direct Walsh evaluation gives
\[
 U={1\over4},\qquad V_1={7\over4},\qquad {V_1\over U}=7.
 \tag{10}
\]
This witness has negative frozen residual, so it refutes only the proposed
unconditional transfer, not (GU-cut) on minimal positive-residual sources.

Second, the adaptive-root single-use theorem bounds root rows encountered
inside one physical read-once restriction tree.  The static sum \(K^{(1)}\)
contains all ordered pairs \((i,j)\) at the common parent.  No source-pinned
lemma currently embeds those mutually incompatible first queries into one
tree without reuse or a terminal remainder.  Hence the adaptive theorem
cannot replace \(V_1/2\) in (GU-cut) as presently stated.

The A099 amplified-extension inequality is a separate sufficient route to
FEI on a critical profile; it does not assert the universal GU inequality
and supplies no algebraic identification of its posterior chain functional
with \(V_1\) or \(D_{\ge2}\).  It is therefore retained as an independent
open route, not inserted as a dependency of (GU-cut).

\subsection*{Exact open leaf}

The leading universal-GU chain now has one explicit leaf: prove (GU-cut)
using (8), or refute that restricted inequality with a genuine Boolean
tensor-prime positive-residual family.  Every root and singleton KL row is
already paid; the unpaid KL sum starts at derivative order two, together
with the explicitly retained reuse functional \(V_1/2\).
